\section{Method}
\label{sec:method}

\subsection{Problem setting}

Let an encoder $E_\theta$ map an image $\mathbf{x} \in \R^{3 \times H
\times W}$ to a feature map $\ze \in \R^{d \times h \times w}$. A
codebook $\CB = \{\mathbf{e}_k\}_{k=1}^{K} \subset \R^{d}$ contains
$K$ vectors. At each spatial location the quantizer performs
nearest-neighbour lookup
\begin{equation}
  k^* = \arg\min_{k \in [K]} \,\|\ze - \mathbf{e}_k\|_2, \qquad
  \qvec := \mathbf{e}_{k^*},
  \label{eq:nnlookup}
\end{equation}
and emits $\zq = \qvec$. A decoder $D_\phi$ reconstructs
$\hat{\mathbf{x}} = D_\phi(\zq)$. Training minimises
\begin{equation}
  \Loss(\theta,\phi,\CB) \;=\;
    \underbrace{\|\mathbf{x} - \hat{\mathbf{x}}\|_1}_{\text{reconstruction}} \;+\;
    \beta\,\underbrace{\|\ze - \sg{\qvec}\|_2^2}_{\text{commitment}} \;+\;
    \text{(EMA codebook update)}.
  \label{eq:loss}
\end{equation}
The EMA update replaces an explicit codebook loss and decays
cluster-membership statistics with rate
$\gamma$~\citep{razavi2019vqvae2}. The difficulty is the $\arg\min$
in~\eqref{eq:nnlookup}: it is piecewise constant and therefore has
zero gradient with respect to $\ze$ almost everywhere.

\subsection{The straight-through estimator, re-examined}

The standard remedy is to substitute the quantizer with the
identity-gradient surrogate
\begin{equation}
  \zq^{\mathrm{ste}} \;=\; \ze + \sg{\qvec - \ze},
  \label{eq:ste}
\end{equation}
so that the forward output is $\qvec$ but the backward derivative is
$\partial \zq^{\mathrm{ste}} / \partial \ze = \mathbf{I}$. The
identity Jacobian has two consequences that motivate this work.
First, \emph{the gradient the encoder receives does not depend on
which code was selected}, because
$\partial \zq^{\mathrm{ste}} / \partial \ze$ does not involve $\qvec$
at all. Second, \emph{the angular relationship between $\ze$ and
$\qvec$ is discarded}. Intuitively, if a particular $\ze$ is mapped
to a $\qvec$ that is nearly orthogonal to it, the encoder should be
pushed strongly to realign its output; under STE it is pushed only by
the reconstruction loss, which propagates through the decoder and
may not expose the misalignment directly.

\subsection{Rotation-plus-rescaling identity}

We now derive an identity that will act as the backward surrogate.
Assume for this subsection that $\ze \neq \mathbf{0}$ and
$\qvec \neq \mathbf{0}$; we handle degenerate cases below. Define
the unit vectors $\hat{\ze} = \ze/\|\ze\|$ and
$\hat{\qvec} = \qvec/\|\qvec\|$ and their difference
$\mathbf{u} = \hat{\ze} - \hat{\qvec}$. When $\hat{\ze} \neq
\hat{\qvec}$ the Householder reflection
\begin{equation}
  \rmat \;=\; \mathbf{I} - 2\,\vvec\vvec^\top,
  \qquad
  \vvec \;=\; \mathbf{u}/\|\mathbf{u}\|,
  \label{eq:householder}
\end{equation}
is a symmetric, orthogonal involution that satisfies
$\rmat\,\hat{\ze} = \hat{\qvec}$. Indeed, writing
$c = \hat{\ze}^\top\hat{\qvec}$ gives $\|\mathbf{u}\|^2 = 2(1 - c)$ and
\begin{equation}
  \rmat\,\hat{\ze}
    = \hat{\ze} - 2(\vvec^\top\hat{\ze})\,\vvec
    = \hat{\ze} - 2\cdot\tfrac{1-c}{\|\mathbf{u}\|}\cdot\vvec
    = \hat{\ze} - \|\mathbf{u}\|\,\vvec
    = \hat{\ze} - \mathbf{u}
    = \hat{\qvec}.
  \label{eq:reflection-check}
\end{equation}
Scaling by $s = \|\qvec\|/\|\ze\|$ and multiplying both sides of
$\rmat\hat{\ze} = \hat{\qvec}$ by $\|\qvec\|$ gives
\begin{equation}
  s\,\rmat\,\ze \;=\; \qvec.
  \label{eq:identity}
\end{equation}
Equation~\eqref{eq:identity} is a \emph{closed-form linear map} from
$\ze$ to $\qvec$: a rotation-reflection and an isotropic scaling, no
free parameters. When $\hat{\ze} = \hat{\qvec}$ (i.e. already aligned)
we take $\rmat = \mathbf{I}$; the scaling $s$ still matches the norms.

\subsection{The \ours{} estimator}

We now use~\eqref{eq:identity} as the backward route. The estimator
is
\begin{equation}
  \boxed{\;\zq^{\,\ours} \;=\; \sg{s\rmat}\,\ze \;+\; \sg{\qvec - \sg{s\rmat}\,\ze}\,.\;}
  \label{eq:our-estimator}
\end{equation}
The first term carries the gradient and has Jacobian $s\rmat$ with
respect to $\ze$ because both $s$ and $\rmat$ are stop-gradients.
The second term is an entirely detached residual that forces the
forward output to equal $\qvec$ bit-exactly. Consequently
\begin{equation}
  \underbrace{\zq^{\,\ours}}_{\text{forward}} \;=\; \qvec,
  \qquad
  \underbrace{\partial \zq^{\,\ours} / \partial \ze}_{\text{backward}} \;=\; s\rmat.
  \label{eq:fwd-bwd}
\end{equation}
The encoder therefore receives gradient $g_{\ze} = (s\rmat)^\top
g_{\zq} = s\rmat g_{\zq}$ (since $\rmat^\top = \rmat$), a
\emph{scaled rotation} of the decoder-side gradient that preserves
its magnitude (up to $s$, which is constant in $\ze$) while
reorienting it through the direction that the forward lookup
implicitly pointed to.

\paragraph{Ablation decomposition.}
Equation~\eqref{eq:our-estimator} admits three natural ablations that
we use in Section~\ref{sec:results}. Setting $\rmat = \mathbf{I}$
while keeping $s$ (``rescale-only'') corresponds to propagating
$s\,\ze$ and discards the angular correction. Setting $s = 1$ while
keeping $\rmat$ (``rotation-only'') propagates $\rmat\,\ze$ and
discards the norm correction. Setting both to identity recovers the
standard STE in~\eqref{eq:ste}. Each variant shares the same forward
pass, so any differences measured in validation reconstruction,
codebook usage, or downstream performance are attributable to the
gradient routing alone.

\subsection{Why Householder?}

The map from $\ze$ to $\qvec$ up to rescaling is determined by an
element of the orthogonal group $O(d)$ only up to the $(d{-}1)$-sphere
of rotations that fix $\hat{\qvec}$. We choose the specific element
that lies in the two-dimensional plane spanned by $\hat{\ze}$ and
$\hat{\qvec}$; this is the Householder reflection
in~\eqref{eq:householder} and is the unique rotation whose action
on vectors orthogonal to that plane is the identity. Three practical
consequences follow.
(i) $\rmat$ is rank-one-plus-identity and requires only the vector
$\vvec \in \R^{d}$ to describe; the action
$\rmat\mathbf{x} = \mathbf{x} - 2(\vvec^\top\mathbf{x})\vvec$ costs
$\mathcal{O}(d)$ per token in both memory and compute, strictly
cheaper than an extra $d \times d$ linear layer.
(ii) Because $\rmat$ is an involution, $\rmat^2 = \mathbf{I}$, which
lets us verify the identity in~\eqref{eq:reflection-check} numerically
to machine precision.
(iii) Householder reflections are the building blocks of stable
$QR$ decompositions; our operator inherits the standard numerical
recipes for handling small $\|\mathbf{u}\|$ (see below).

\subsection{Numerical stability}

The naive implementation of~\eqref{eq:householder} divides by
$\|\mathbf{u}\| = \sqrt{2(1-c)}$, which is small when $\ze$ and
$\qvec$ are nearly parallel. In that regime $\hat{\ze} \approx
\hat{\qvec}$, so both $\rmat \approx \mathbf{I}$ and $s \approx 1$ and
the surrogate already behaves like the STE up to the norm
correction. Rather than branch, we use the stop-gradient decomposition
in~\eqref{eq:our-estimator}: the second term forces the forward
output to be $\qvec$ regardless of how ill-conditioned $\rmat$ is,
and the backward gradient is dominated by $\mathbf{I}$ whenever
$\rmat$ is numerically close to $\mathbf{I}$. We guard
$\|\mathbf{u}\|$ against division-by-zero with a small
$\varepsilon$ (set to $10^{-6}$ in practice) and zero out the
reflection component whenever $\|\mathbf{u}\| < \varepsilon$, in
which case $\zq^{\,\ours}$ reduces smoothly to $s\,\ze + \sg{\qvec - s\,\ze}$.
This behaviour is verified empirically in Section~\ref{sec:results}:
our unit tests confirm that the autograd Jacobian agrees with
$s\rmat$ to within $10^{-7}$ in single precision when $\ze$ is well
separated from $\qvec$, and that the forward output matches $\qvec$ to
within $10^{-5}$ for inputs drawn after codebook warm-up.

\subsection{Algorithm and cost}

Algorithm~\ref{alg:rotquant} summarises the per-token forward pass.
All operations are $\mathcal{O}(d)$ and trivially parallelise over
the batch, spatial grid, and codebook dimensions; the operator is
implemented with three fused elementwise tensor ops and no matrix
materialisation. Memory overhead is $d$ floats per token for
$\vvec$ (freed after the backward pass) and one scalar for $s$. When
compared against a baseline VQ-VAE trained under STE, the forward
and backward costs of~\eqref{eq:our-estimator} differ by less than a
constant factor independent of the codebook size $K$, so throughput
is preserved even for codebooks of $16{,}384$ entries
(Section~\ref{sec:results}, compute-overhead experiment).

\begin{algorithm}[t]
\caption{\ours{} forward pass (per latent token).}
\label{alg:rotquant}
\begin{algorithmic}[1]
  \Require encoder output $\ze \in \R^{d}$, codebook $\{\mathbf{e}_k\} \subset \R^{d}$, $\varepsilon > 0$
  \State $k^* \gets \arg\min_k \|\ze - \mathbf{e}_k\|_2$;\quad $\qvec \gets \mathbf{e}_{k^*}$
  \State $\hat{\ze} \gets \ze/\max(\|\ze\|,\varepsilon)$;\quad $\hat{\qvec} \gets \qvec/\max(\|\qvec\|,\varepsilon)$
  \State $\mathbf{u} \gets \hat{\ze} - \hat{\qvec}$;\quad $\vvec \gets \mathbf{u}/\max(\|\mathbf{u}\|,\varepsilon)$
  \State $s \gets \sg{\|\qvec\|/\|\ze\|}$
  \State carrier $\gets s\,\big(\ze - 2(\vvec^\top \ze)\,\vvec\big)$ \Comment{backward Jacobian $= s\rmat$ (with $\vvec$ detached)}
  \State $\zq \gets \text{carrier} + \sg{\qvec - \text{carrier}}$ \Comment{forward $= \qvec$ bit-exact}
  \State \Return $\zq$
\end{algorithmic}
\end{algorithm}

\subsection{Relationship to prior gradient estimators}

Our estimator can be placed within the taxonomy of prior work as
follows. \citet{bengio2013ste} and \citet{vandenoord2017vqvae} use
$\partial \zq / \partial \ze = \mathbf{I}$, a zeroth-order
approximation. \citet{huh2023straighteningste} introduce an
alternating-minimisation scheme and an additional rotation-like
residual to smooth optimisation; their backward Jacobian is no
longer the identity but is not derived from the closed-form
rotation~\eqref{eq:identity}. \citet{shah2024decoupledste} decouple
forward and backward identities but retain an identity-like
Jacobian. The Gumbel-Softmax~\citep{jang2017gumbel} and
concrete~\citep{maddison2017concrete} estimators replace the
$\arg\min$ with a softmax-weighted combination of all codes; their
forward output is a soft mixture and the backward gradient depends on
the full code distribution. FSQ~\citep{mentzer2024fsq} eliminates the
codebook and uses a per-channel STE on a bounded-and-rounded
projection. NSVQ~\citep{vali2022nsvq} replaces the residual with
noise. Closest is the concurrent work of~\citet{fifty2024rotation},
which independently arrives at a rotation-based backward. Our
treatment differs in deriving the rotation as a closed-form
Householder operator, in ablating its components against scalar
rescaling, and in analysing the degenerate
$\hat{\ze} \approx \hat{\qvec}$ regime where the reflection
matrix is ill-conditioned.
